% GENERATED FILE. Edit canonical lessons, then run npm run generate:books.
% canonical-source-hash: fnv1a64:e3ec6e74dc1a6dca
% canonical-lessons: JA-C437-kibishii, JA-C437-reisei, JA-C437-seijitsu, JA-C437-kenkyo, JA-C437-ijiwaru, JA-R437-review-432-437

\chapter{Strict, Sincere or Mean}
\label{ch:ja-strict-sincere-or-mean}
\hlchaptermodality{437}

\begin{chapteropening}
\textbf{By the end of this chapter:} \emph{I can describe a strict, sincere or mean person.}

Complete the last lesson of chapter 437: I can describe a strict, sincere or mean person.
\end{chapteropening}

\section[\texorpdfstring{\ja{きびしい} (kibishii)}{きびしい (kibishii)}]{\ja{きびしい} (kibishii) — strict; severe}

\label{lesson:JA-C437-kibishii}



\begin{quote}
Before the new one: say the Japanese for shy with strangers, then the Japanese for to behave.
\end{quote}

\subsection*{You'll want to know: \ja{きびしい}}
\textbf{\ja{きびしい}} — \emph{kibishii} — \textquotedblleft{}strict; severe\textquotedblright{}.

\begin{quote}
\textbf{\ja{じょうしは} \ja{しごとに} \ja{きびしいです。}} — \emph{jōshi wa shigoto ni kibishii desu} — \textquotedblleft{}My boss is strict about work.\textquotedblright{}
\end{quote}

An i-adjective.

\begin{cousinweb}
\textbf{\ja{きびしい}} (\emph{kibishii}) is a native Japanese word (wago), inherited rather than borrowed from Chinese; its kanji spelling was attached to it later (JMdict).
\end{cousinweb}

\begin{grammarlens}[title={What you've built}]
One more word for describing what you lived and felt.
\end{grammarlens}

\subsection*{Your turn}
\begin{itemize}
\raggedright
  \item \emph{Say it:} \emph{kibishii}
  \item \emph{Say it:} \emph{kibishii}, once more
  \item \emph{Say it:} \emph{jōshi wa shigoto ni kibishii desu} — My boss is strict about work.
  \item \emph{Recall:} say the Japanese for shy with strangers, then the Japanese for to behave, then say \emph{kibishii} again
\end{itemize}

\begin{tcolorbox}[breakable,colback=teal!4,colframe=teal!35!black,title={Before you move on}]
What does \ja{きびしい} mean? (\textquotedblleft{}Strict; severe\textquotedblright{}.) Say it once more.
\end{tcolorbox}
\section[\texorpdfstring{\ja{れいせい} (reisei)}{れいせい (reisei)}]{\ja{れいせい} (reisei) — composed; cool-headed}

\label{lesson:JA-C437-reisei}



\begin{quote}
Before the new one: say the Japanese for to behave, then the Japanese for strict.
\end{quote}

\subsection*{You'll want to know: \ja{れいせい}}
\textbf{\ja{れいせい}} — \emph{reisei} — \textquotedblleft{}composed; cool-headed\textquotedblright{}.

\begin{quote}
\textbf{\ja{かれは} \ja{いつも} \ja{れいせいです。}} — \emph{kare wa itsumo reisei desu} — \textquotedblleft{}He keeps a cool head at all times.\textquotedblright{}
\end{quote}

A na-adjective: \ja{な} before a noun.

\begin{cousinweb}
\textbf{\ja{れいせい}} (\emph{reisei}) is Sino-Japanese (kango): two Chinese characters read with their on'yomi, readings borrowed from Chinese. They are \emph{rei}, \textquotedblleft{}cool\textquotedblright{} (Middle Chinese \emph{lengX}) and \emph{sei}, \textquotedblleft{}quiet\textquotedblright{} (\emph{dzjengX}) (KANJIDIC2; Middle Chinese in Baxter's transcription).
\end{cousinweb}

\begin{grammarlens}[title={What you've built}]
One more word for describing what you lived and felt.
\end{grammarlens}

\subsection*{Your turn}
\begin{itemize}
\raggedright
  \item \emph{Say it:} \emph{reisei}
  \item \emph{Say it:} \emph{reisei}, once more
  \item \emph{Say it:} \emph{kare wa itsumo reisei desu} — He keeps a cool head at all times.
  \item \emph{Recall:} say the Japanese for to behave, then the Japanese for strict, then say \emph{reisei} again
\end{itemize}

\begin{tcolorbox}[breakable,colback=teal!4,colframe=teal!35!black,title={Before you move on}]
What does \ja{れいせい} mean? (\textquotedblleft{}Composed; cool-headed\textquotedblright{}.) Say it once more.
\end{tcolorbox}
\section[\texorpdfstring{\ja{せいじつ} (seijitsu)}{せいじつ (seijitsu)}]{\ja{せいじつ} (seijitsu) — sincere; faithful}

\label{lesson:JA-C437-seijitsu}



\begin{quote}
Before the new one: say the Japanese for strict, then the Japanese for composed.
\end{quote}

\subsection*{You'll want to know: \ja{せいじつ}}
\textbf{\ja{せいじつ}} — \emph{seijitsu} — \textquotedblleft{}sincere; faithful\textquotedblright{}.

\begin{quote}
\textbf{\ja{じょうしは} \ja{せいじつです。}} — \emph{jōshi wa seijitsu desu} — \textquotedblleft{}My boss is sincere.\textquotedblright{}
\end{quote}

A na-adjective: \ja{な} before a noun.

\begin{cousinweb}
\textbf{\ja{せいじつ}} (\emph{seijitsu}) is Sino-Japanese (kango): two Chinese characters read with their on'yomi, readings borrowed from Chinese. They are \emph{sei}, \textquotedblleft{}sincerity\textquotedblright{} (Middle Chinese \emph{dzyeng}) and \emph{jitsu}, \textquotedblleft{}reality\textquotedblright{} (\emph{zyit}) (KANJIDIC2; Middle Chinese in Baxter's transcription).
\end{cousinweb}

\begin{grammarlens}[title={What you've built}]
One more word for describing what you lived and felt.
\end{grammarlens}

\subsection*{Your turn}
\begin{itemize}
\raggedright
  \item \emph{Say it:} \emph{seijitsu}
  \item \emph{Say it:} \emph{seijitsu}, once more
  \item \emph{Say it:} \emph{jōshi wa seijitsu desu} — My boss is sincere.
  \item \emph{Recall:} say the Japanese for strict, then the Japanese for composed, then say \emph{seijitsu} again
\end{itemize}

\begin{tcolorbox}[breakable,colback=teal!4,colframe=teal!35!black,title={Before you move on}]
What does \ja{せいじつ} mean? (\textquotedblleft{}Sincere; faithful\textquotedblright{}.) Say it once more.
\end{tcolorbox}
\section[\texorpdfstring{\ja{けんきょ} (kenkyo)}{けんきょ (kenkyo)}]{\ja{けんきょ} (kenkyo) — modest; humble}

\label{lesson:JA-C437-kenkyo}



\begin{quote}
Before the new one: say the Japanese for composed, then the Japanese for sincere.
\end{quote}

\subsection*{You'll want to know: \ja{けんきょ}}
\textbf{\ja{けんきょ}} — \emph{kenkyo} — \textquotedblleft{}modest; humble\textquotedblright{}.

\begin{quote}
\textbf{\ja{あの} \ja{せんしゅは} \ja{けんきょです。}} — \emph{ano senshu wa kenkyo desu} — \textquotedblleft{}That athlete is modest.\textquotedblright{}
\end{quote}

A na-adjective: \ja{な} before a noun.

\begin{cousinweb}
\textbf{\ja{けんきょ}} (\emph{kenkyo}) is Sino-Japanese (kango): two Chinese characters read with their on'yomi, readings borrowed from Chinese. They are \emph{ken}, \textquotedblleft{}be modest\textquotedblright{} (Middle Chinese \emph{khem}) and \emph{kyo}, \textquotedblleft{}emptiness\textquotedblright{} (\emph{xjo}) (KANJIDIC2; Middle Chinese in Baxter's transcription).
\end{cousinweb}

\begin{grammarlens}[title={What you've built}]
One more word for describing what you lived and felt.
\end{grammarlens}

\subsection*{Your turn}
\begin{itemize}
\raggedright
  \item \emph{Say it:} \emph{kenkyo}
  \item \emph{Say it:} \emph{kenkyo}, once more
  \item \emph{Say it:} \emph{ano senshu wa kenkyo desu} — That athlete is modest.
  \item \emph{Recall:} say the Japanese for composed, then the Japanese for sincere, then say \emph{kenkyo} again
\end{itemize}

\begin{tcolorbox}[breakable,colback=teal!4,colframe=teal!35!black,title={Before you move on}]
What does \ja{けんきょ} mean? (\textquotedblleft{}Modest; humble\textquotedblright{}.) Say it once more.
\end{tcolorbox}
\section[\texorpdfstring{\ja{いじわる} (ijiwaru)}{いじわる (ijiwaru)}]{\ja{いじわる} (ijiwaru) — spiteful; mean}

\label{lesson:JA-C437-ijiwaru}



\begin{quote}
Before the new one: say the Japanese for sincere, then the Japanese for modest.
\end{quote}

\subsection*{You'll want to know: \ja{いじわる}}
\textbf{\ja{いじわる}} — \emph{ijiwaru} — \textquotedblleft{}spiteful; mean\textquotedblright{}.

\begin{quote}
\textbf{\ja{いもうとに} \ja{いじわるを} \ja{しないで} \ja{ください。}} — \emph{imōto ni ijiwaru o shinaide kudasai} — \textquotedblleft{}Please do not be mean to your younger sister.\textquotedblright{}
\end{quote}

A na-adjective: \ja{な} before a noun.

\begin{cousinweb}
\textbf{\ja{いじわる}} (\emph{ijiwaru}) is a mixed word: Sino-Japanese \emph{iji}, \textquotedblleft{}disposition\textquotedblright{}, from \emph{i}, \textquotedblleft{}mind\textquotedblright{} (Middle Chinese \emph{'iH}), and \emph{ji}, \textquotedblleft{}ground\textquotedblright{} (\emph{dijH}), borrowed from Chinese, joined to native \emph{waru}, \textquotedblleft{}bad\textquotedblright{}, the stem of \emph{warui} (JMdict, KANJIDIC2; Middle Chinese in Baxter's transcription).
\end{cousinweb}

\begin{grammarlens}[title={What you've built}]
One more word for describing what you lived and felt.
\end{grammarlens}

\subsection*{Your turn}
\begin{itemize}
\raggedright
  \item \emph{Say it:} \emph{ijiwaru}
  \item \emph{Say it:} \emph{ijiwaru}, once more
  \item \emph{Say it:} \emph{imōto ni ijiwaru o shinaide kudasai} — Please do not be mean to your younger sister.
  \item \emph{Recall:} say the Japanese for sincere, then the Japanese for modest, then say \emph{ijiwaru} again
\end{itemize}

\begin{tcolorbox}[breakable,colback=teal!4,colframe=teal!35!black,title={Before you move on}]
What does \ja{いじわる} mean? (\textquotedblleft{}Spiteful; mean\textquotedblright{}.) Say it once more.
\end{tcolorbox}
\section[(dialogue)]{Looking back: from Worry and Regret to Strict, Sincere or Mean}

\label{lesson:JA-R437-review-432-437}



\begin{quote}
Say the Japanese for to be troubled and to brood.
\end{quote}

\subsection*{Your turn}
Say these aloud:

\begin{itemize}
\raggedright
  \item to be troubled, to brood, to regret, futile, to give up
  \item to yearn after, missed dearly, envious, an interest, admiration
  \item to be grateful, to feel sympathy, pitiable, miserable, to encourage
  \item satisfied, discontented, bored, to smile gently, to fall silent
  \item cheerful, bright, taciturn, shy with strangers, to behave
  \item strict, composed, sincere, modest, spiteful
\end{itemize}

\begin{tcolorbox}[breakable,colback=teal!4,colframe=teal!35!black,title={Before you move on}]
To be troubled? (\textbf{\ja{なやむ}}, \emph{nayamu}.) To brood? (\textbf{\ja{くよくよする}}, \emph{kuyokuyo suru}.) To regret? (\textbf{\ja{こうかいする}}, \emph{kōkai suru}.) Futile? (\textbf{\ja{むなしい}}, \emph{munashii}.) To give up? (\textbf{\ja{あきらめる}}, \emph{akirameru}.) To yearn after? (\textbf{\ja{あこがれる}}, \emph{akogareru}.) Missed dearly? (\textbf{\ja{こいしい}}, \emph{koishii}.) Envious? (\textbf{\ja{うらやましい}}, \emph{urayamashii}.) An interest? (\textbf{\ja{きょうみ}}, \emph{kyōmi}.) Admiration? (\textbf{\ja{かんしん}}, \emph{kanshin}.) To be grateful? (\textbf{\ja{かんしゃする}}, \emph{kansha suru}.) To feel sympathy? (\textbf{\ja{どうじょうする}}, \emph{dōjō suru}.) Pitiable? (\textbf{\ja{きのどく}}, \emph{kinodoku}.) Miserable? (\textbf{\ja{みじめ}}, \emph{mijime}.) To encourage? (\textbf{\ja{はげます}}, \emph{hagemasu}.) Satisfied? (\textbf{\ja{まんぞく}}, \emph{manzoku}.) Discontented? (\textbf{\ja{ふまん}}, \emph{fuman}.) Bored? (\textbf{\ja{たいくつ}}, \emph{taikutsu}.) To smile gently? (\textbf{\ja{ほほえむ}}, \emph{hohoemu}.) To fall silent? (\textbf{\ja{だまる}}, \emph{damaru}.) Cheerful? (\textbf{\ja{ほがらか}}, \emph{hogaraka}.) Bright? (\textbf{\ja{あかるい}}, \emph{akarui}.) Taciturn? (\textbf{\ja{むくち}}, \emph{mukuchi}.) Shy with strangers? (\textbf{\ja{ひとみしり}}, \emph{hitomishiri}.) To behave? (\textbf{\ja{ふるまう}}, \emph{furumau}.) Strict? (\textbf{\ja{きびしい}}, \emph{kibishii}.) Composed? (\textbf{\ja{れいせい}}, \emph{reisei}.) Sincere? (\textbf{\ja{せいじつ}}, \emph{seijitsu}.) Modest? (\textbf{\ja{けんきょ}}, \emph{kenkyo}.) Spiteful? (\textbf{\ja{いじわる}}, \emph{ijiwaru}.)
\end{tcolorbox}
